\documentclass[11pt,letterpaper]{article}
\usepackage[T1]{fontenc}
\usepackage{lmodern}
\usepackage[margin=1in,headheight=15pt]{geometry}
\usepackage{amsmath,amssymb,amsthm,mathtools}
\usepackage{booktabs,longtable,tabularx,array}
\usepackage{enumitem,microtype,xcolor,fancyhdr}
\usepackage{hyperref}
\definecolor{inkblue}{HTML}{164C63}
\hypersetup{colorlinks=true,linkcolor=inkblue,citecolor=inkblue,urlcolor=inkblue,
  pdftitle={Dyadic descaling and delayed Pell sign recovery for Diophantine consistency},
  pdfauthor={OpenAI-assisted research prepared for Vladimir Reshetnikov}}
\newcommand{\Z}{\mathbb Z}
\newcommand{\N}{\mathbb N}
\newcommand{\Q}{\mathbb Q}
\newcommand{\Pos}{\mathbb Z_{>0}}
\newcommand{\Con}{\operatorname{Con}}
\newcommand{\Refut}{\operatorname{Ref}}
\newcommand{\ZFC}{\mathsf{ZFC}}
\newcommand{\ZFp}{\mathsf{ZF}}
\newcommand{\vTwo}{\nu_2}
\newcommand{\sha}{\texttt{ec80841c07e91be8956ce6ed4af9a9c0b782d3d9}}
\newtheorem{theorem}{Theorem}[section]
\newtheorem{lemma}[theorem]{Lemma}
\newtheorem{proposition}[theorem]{Proposition}
\newtheorem{corollary}[theorem]{Corollary}
\theoremstyle{definition}
\newtheorem{definition}[theorem]{Definition}
\theoremstyle{remark}
\newtheorem{remark}[theorem]{Remark}
\numberwithin{equation}{section}
\setlist{itemsep=3pt,topsep=5pt}
\setlength{\emergencystretch}{2em}
\pagestyle{fancy}
\fancyhf{}
\lhead{\small Diophantine consistency: two reductions}
\rhead{\small ProveIt research package}
\cfoot{\thepage}
\renewcommand{\headrulewidth}{0.3pt}
\title{\textbf{Dyadic descaling and delayed Pell sign recovery}\\[4pt]
  \large Two explicit reductions for Diophantine expressions of $\Con(\ZFC)$}
\author{Prepared for Vladimir Reshetnikov\\
  \small Research, implementation, and separate proof review with ChatGPT}
\date{23 October 2026}

\begin{document}
\maketitle
\thispagestyle{empty}
\begin{abstract}
We construct two explicit polynomial equations whose unsolvability is
equivalent to $\Con(\ZFC)$, using precisely stated calibration theorems
already present in ProveIt. The first equation has ten strictly positive
unknowns, exact total degree $706$, and a division-free circuit with $133$
binary arithmetic operations. Its immediate parent has degree $816$ and
$124$ operations. The essential change is to separate four rational square
classes by distinct powers of two. This removes a strong-Pell scale factor
that had been needed for real-size separation. Three centers and a quadratic
transport majorant then determine the lower exact degree.

The second equation is a quartic in $780$ unrestricted integers, defined
both by $434$ explicit quadratic residuals and by all $9,628$ collected
monomials. It improves the $888$-integer parent in the latest direct
certificate package. Its largest absolute coefficient is
$2,445,790,974,704$, unchanged from that parent. This route uses the direct
C222 proof-certificate calibration rather than K215 or Jones's universal
equation. The reduction removes $36$ nonnegativity guards by proving their
consequences in a strict order: elementary signs, the Pell frame, indexed
powers, and finally signs of power quotients.

The new algebraic and number-theoretic reductions are proved here;
the two source calibrations are inherited, not silently replaced by
experiments. Literal data, exact rebuilding programs, source hashes,
audit receipts, and an independent AI-assisted proof review accompany
the article. We do not claim global minimality, a new universal-variable
record, or an end-to-end proof-assistant certificate.
\end{abstract}
\clearpage
\tableofcontents
\clearpage

\section{The target and the two results}
\label{sec:overview}

Fix the usual recursive presentation of ZFC used by the repository's
calibration. Write $\Refut_{\ZFC}(p)$ for its primitive-recursive predicate
that $p$ is a finite syntactic refutation, and abbreviate
\[
  \Con(\ZFC)\;\equiv\;\neg\exists p\in\N\;\Refut_{\ZFC}(p).
\]
The desired assertion for a polynomial $P$ and an explicitly specified
domain $D$ is
\begin{equation}\label{eq:target}
 \ZFC\vdash
 \left(\forall \boldsymbol x\in D^v\;P(\boldsymbol x)\ne0\right)
 \longleftrightarrow \Con(\ZFC).
\end{equation}
The useful way to establish this is to prove, internally by ordinary
arithmetizable mathematics, that a zero exists exactly when a finite
refutation exists. There is no assertion here that ZFC proves its own
consistency, or that a numerical search could establish unsolvability.

The word ``simplest'' requires a choice of cost. Unknown count, degree,
arithmetic circuit length, number of monomials, coefficient height,
compressed coefficient storage, and proof dependencies are different
measures. A sparse quartic with hundreds of unknowns can have much smaller
coefficients than a compact formula in ten unknowns. This work therefore
states concrete improvements and their costs rather than an unproved
absolute optimum.

\begin{table}[ht]
\centering\small
\begin{tabular}{@{}lrrrr@{}}
\toprule
Construction & Unknowns & Domain & Degree & Further size measure\\
\midrule
Immediate scaled-norm parent & 10 & $\Pos$ & 816 & 124 operations\\
Dyadic construction here & 10 & $\Pos$ & \textbf{706} & 133 operations\\
Direct parent & 888 & $\Z$ & 4 & 10,156 monomials\\
Delayed-sign construction here & \textbf{780} & $\Z$ & 4 & 9,628 monomials\\
\bottomrule
\end{tabular}
\caption{Exact comparisons to the two pinned immediate parents.
Domains and coefficient bindings are kept distinct.}
\label{tab:main}
\end{table}

The first change reduces degree by $110/816\approx13.48\%$ and adds nine
arithmetic operations. It does not dominate its parent in every cost.
The second removes $108/888\approx12.16\%$ of the integer unknowns while
preserving degree and coefficient height. It is a reduction within the
direct proof-certificate architecture, not a claim that no smaller
integer quartic occurs elsewhere in the repository.

All comparisons use the fixed repository snapshot\newline
\sha. The reviewed
collection already contains nine-positive-variable and ten-integer
equations, and earlier native-calibration quartics have far fewer than
$780$ integer unknowns. The latter use a different coefficient regime.
The direct $407$-natural-variable parent is also retained unchanged by
our new integer-domain conversion. These comparisons prevent a domain
change or an obsolete baseline from being mistaken for a record.

\subsection{What changes mathematically}

The ten-positive-variable parent makes several real radicals grow at
different rates, so an irrational square class cannot disappear from a
signed sum. That requirement inflates its strong Pell ordinate.
Here we arrange instead that each radicand is a unit at the prime $2$,
then weight the four roots by $1,2,4,8$. Within any common rational square
class, the weighted coefficients have distinct powers of two. The least
power cannot cancel. The same rational-root conclusion follows without
size ordering, and one scale factor can be deleted.

The direct construction addresses a different source of growth.
An integer version of a natural-number certificate often attaches three
new variables to each sign requirement. Some sign requirements become
theorems after the arithmetic meaning of the certificate is established.
We prove those meanings before inferring the remaining signs, and record
the exact dependency order. Removing $36$ such requirements saves
$3\cdot36=108$ integer variables.

Neither the use of valuations nor the elementary quotient recurrence is
claimed as a new general principle. The contributions are their concrete
applications, the compatible proof ordering, the complete fixed
polynomials, and the quantified improvements at the pinned source.
General work on universal pairs illustrates why degree and variable
count must be tracked separately \cite{bayer2026}; Sun's ten-integer
theorem concerns a different, universal reduction \cite{sun2026}.
Neither new endpoint here invokes Sun's ten-integer theorem.

\section{Inherited calibration and precise proof boundaries}
\label{sec:calibration}

\subsection{The complete fixed coefficient tuple}

For the first construction the fixed numerals are
\[
   \beta=2^{590069}-1,\qquad K_0,\qquad M_C,\qquad M_F.
\]
The last three are literal integers supplied by the unchanged compact44
coefficient files. Each file is a gzip-compressed little-endian binary
integer. Their complete compressed and uncompressed SHA-256 bindings
appear in the circuit JSON and its builder. The designation $M_F$ means
the \emph{already shifted} high mask. There is no additional shift.
These fixed numerals satisfy
\begin{equation}\label{eq:ports}
 K_0>0,\qquad 0<M_C<\beta<M_F<2\beta.
\end{equation}
They are part of the polynomial's definition, not existential parameters,
an unspecified enumeration index, or a truth-dependent choice.
The four compressed payloads together occupy about nine kilobytes;
their uncompressed representation has $295,388$ bytes.

\begin{table}[ht]
\centering\small
\begin{tabular}{@{}llr@{}}
\toprule
Numeral & Supplied file & Bit length\\
\midrule
$\beta$ & \texttt{beta.bin.gz} & 590,069\\
$K_0$ & \texttt{K.bin.gz} & 592,884\\
$M_C$ & \texttt{MC.bin.gz} & 590,069\\
$M_F$ & \texttt{MF\_source.bin.gz} & 590,070\\
\bottomrule
\end{tabular}
\caption{The whole compact44 binding. No ports from another tuple are mixed in.}
\end{table}

\subsection{The native positive source interface}

Here is the exact inherited theorem used for the ten-positive equation.
It is Theorem \texttt{thm:source} of package 61, including the canonical
promises immediately following it \cite{source61}. Define, for positive
source coordinates
\[
 J,F,\alpha,z,f,h,i,T_{\rm old},s,w,\tau,\eta,\zeta,y,\gamma,
\]
the registers
\begin{align*}
 q&=\beta J+1,& C&=q-F-\alpha,& X&=wq,& Y&=sq^3,\\
 E&=XY,& k&=\eta+\zeta,& c&=kY+\eta,& a&=E+Y,\\
 H_0&=4a+3,& \Delta_0&=(a+2)^2-1,& \ell&=EkY,& I_0&=ic^2.
\end{align*}
Set
\begin{align*}
 R&=((q-1)(q-F)+\alpha)(q^2-1)+(M_C+qM_F)J,\\
 D&=X+ac+H_0\gamma,\qquad
 V_{\rm old}=c(T_{\rm old}f-1)-Rf^2.
\end{align*}
The source equations are
\begin{align}
 \tau^2-\ell(\ell+k)&=1,&D^2-\Delta_0c^2&=1,\nonumber\\
 f^2-\Delta_0I_0^2&=1,&
 \Delta_0(f^2-1)(V_{\rm old}^2-y^2)+y^2&=1,\label{eq:source}\\
 k-hE-R&=1,&(K_0+w)C+(q-F)-z(q-1)&=1.\nonumber
\end{align}

\begin{theorem}[Inherited calibrated positive source]\label{thm:calibration}
For the entire tuple just specified, a positive solution of
\eqref{eq:source} yields a finite syntactic refutation of ZFC.
Conversely, a finite refutation yields positive solutions with
\begin{align*}
 q&=B^N,\quad B=\beta+1,\quad N>1\text{ odd},\quad
 J=1+B+\cdots+B^{N-1},\\
 R&\equiv3\pmod4,\qquad X=2^R,\qquad
 c=\psi_{a+2}(R),\\
 k&=2\psi_{2XY^2+1}((R+1)/2),\qquad
 \tau=\chi_{2XY^2+1}((R+1)/2).
\end{align*}
The two implications are ordinary mathematics formalizable in ZFC.
\end{theorem}

Here $\chi_A(n)+\psi_A(n)\sqrt{A^2-1}
=(A+\sqrt{A^2-1})^n$. The source theorem composes the native compiler,
its compact44 adapter, and the literal blank-K215/refutation calibration.
We inherit that composition. The source is retained and pinned in the
package; this article does not re-prove every compiler and machine lemma.
The immediate degree-$816$ proof \cite{parent816} gives the precise Pell
rank, index, parity, ratio, and auxiliary-completeness interfaces used
below, and is included unchanged for inspection.

\subsection{The independent direct source interface}

The second route starts with the exact $690$-natural-variable quadratic
residual system in the pinned incoming package
\emph{ConZFC\_407\_888\_Quartics} \cite{direct407}. Its natural zero set
is calibrated to finite ZFC refutations through the direct C222/Horn
certificate, beta coding, binomial arithmetic, and shared Pell frames.
Its existing positive-to-integer conversion has $888$ unknowns.
Our transformation changes the sign certificates and then performs exact
monic eliminations. It preserves the source system's natural zero set.
The original archive, including the source explanation and provenance,
is retained byte for byte.

Thus the logical dependency has two independently specified forms:
\[
\begin{split}
 N_{706}=0\text{ in }\Pos^{10}
 &\longleftrightarrow\text{native positive source zero}
 \longleftrightarrow\exists p\;\Refut_{\ZFC}(p),\\
 P_{780}=0\text{ in }\Z^{780}
 &\longleftrightarrow\text{direct natural source zero}
 \longleftrightarrow\exists p\;\Refut_{\ZFC}(p).
\end{split}
\]
Negating each chain gives \eqref{eq:target}.
The new reductions do not upgrade the inherited chain to an exported
object-language ZFC proof or a complete Lean/Rocq formalization.

\section{A degree-706 equation in ten positive integers}
\label{dy:section}

The first construction preserves the ten positive unknowns of the
scaled-norm candidate in the pinned repository, while reducing its exact
total degree from $816$ to $706$.  Its division-free circuit uses $133$
binary arithmetic operations, compared with $124$ for that parent.
This is a degree--operation tradeoff, rather than a claim that one measure
of simplicity captures every consideration.  The coefficient tuple and
its calibration remain those described in Section~\ref{sec:calibration}.
The immediate mathematical parent is~\cite{parent816}; the source interface
and the complete fixed coefficient binding are~\cite{source61,source44}.

The main change is to make all four rational radicands units at the prime
$2$.  Distinct powers of $2$ then separate their rational square classes.
This removes a size-order condition on the strong Pell root and permits
one factor to be removed from its ordinate.  Three centers in the carrier
and a quadratic transport majorant account for the remaining reduction.
We prove the new algebra and positivity arguments explicitly.  The
inherited compiler/calibration and the stated Pell index lemmas retain
precisely the status recorded in Section~\ref{sec:calibration}.

\subsection{The polynomial and its fixed coefficients}
\label{dy:definition}

Write the four fixed integer coefficients as
\[
 \beta=2^{590069}-1,\qquad K_0,\qquad M_C,\qquad M_F.
\]
The complete binary payloads are supplied with the circuit and authenticated
by their compressed and uncompressed SHA-256 digests.  Their verified
bounds are
\begin{equation}
  K_0>0,\qquad 0<M_C<\beta<M_F<2\beta.
  \label{dy:ports}
\end{equation}
Here $M_F$ is already the shifted high mask; it must not be shifted a
second time.  None of these four constants is a quantified unknown.

The ten unknowns are
\begin{equation}
 (j,F,\alpha,i,T,w,Y,n,\eta,\zeta)\in\Pos^{10},
 \qquad \Pos=\{1,2,3,\ldots\}.
 \label{dy:variables}
\end{equation}
Every register in the following definition is an integer polynomial in
these unknowns and the four fixed constants.  Define the outer registers
by
\begin{align}
 J&=2j+1,& q&=\beta J+1,& C&=q-F-\alpha, \notag\\
 X&=wq,& a&=Y(X+1),& k&=\eta+\zeta, \notag\\
 c&=kY+\eta,& \mu&=q^2+1,& \lambda&=2\mu.
 \label{dy:outer}
\end{align}
The scaled main and first Pell parameters are
\begin{align}
 A&=\lambda(a+2)+1,& H&=4A-5,& \Delta&=A^2-1, \notag\\
 a_-&=A-2,& V_0&=X\lambda^2Y^2,&
 D_\tau&=V_0(V_0+1)k^2+1.
 \label{dy:parameters}
\end{align}
The packed index, its excess and the transport numerator are
\begin{align}
 G&=(q-1)(q-F)+\alpha, \notag\\
 R&=G(q^2-1)+(M_C+qM_F)J, \notag\\
 e&=k-R-1, & B_t&=(K_0+w)C-F.
 \label{dy:index}
\end{align}
For the strong and auxiliary Pell registers put
\begin{align}
 I&=ic^2,& D_f&=1+\Delta I^2,& Q&=\Delta^2I^2,
       &\delta&=Q-1, \notag\\
 T'&=T+\alpha,& b&=cT'-R,& V&=D_fb-c, \notag\\
 D_y&=QV^2-1,& D_m&=\Delta c^2+1.
 \label{dy:auxiliary}
\end{align}
In particular, $I$ has no factor $\mu$ other than any accidental
divisibility of its value.  Define the centers and transport weight
\begin{align}
 C_f&=AI,& C_\tau&=V_0k, \notag\\
 L&=(k+F+1)(T'+6),& W&=1+(q-1)L,
 \label{dy:weight}
\end{align}
and the denominator-clearing registers
\begin{align}
 P_1&=q^3(q-1),& M&=P_1\mu^2H, \notag\\
 Z&=P_1\mu^2\bigl[H(n+V+2C_f+4C_\tau)+8(X+a_-c)\bigr]
       \notag\\
  &\hspace{1.2em}+H\bigl[\mu^2((q-1)Y-q^3WB_t)-P_1e\bigr].
 \label{dy:carrier}
\end{align}
For the final norm use the five polynomial arguments
\begin{align}
 s&=\delta Z^2,&
 u_1&=16\delta M^2D_\tau,&
 u_2&=4\delta M^2D_f, \notag\\
 u_3&=M^2D_y,&
 u_4&=64\delta(P_1\mu^2)^2D_m.
 \label{dy:normargs}
\end{align}

Here is a compact fixed norm polynomial in five formal arguments.  Its
letters $a,b,c,d$ are local formal arguments, independent of the similarly
named registers above:
\begin{align}
 \sigma_1&=a+b+c,& \sigma_2&=ab+(a+b)c,& \sigma_3&=abc, \notag\\
 U&=s+d,& P&=U-\sigma_1,& E&=P^2+4(sd-\sigma_2), \notag\\
 G_*&=E^2+16(sdP^2-4\sigma_3U),&
 O_*&=EP-16\sigma_3, \notag\\
 \mathcal F_{16}(s;a,b,c,d)&=G_*^2-64sdO_*^2.
 \label{dy:f16}
\end{align}
The equation delivered here is
\begin{equation}
 \boxed{\;N=\mathcal F_{16}(s;u_1,u_2,u_3,u_4)=0,\quad
         (j,F,\alpha,i,T,w,Y,n,\eta,\zeta)\in\Pos^{10}.\;}
 \label{dy:equation}
\end{equation}
The named registers are abbreviations for substitution, not additional
unknowns.  Equations~\eqref{dy:outer}--\eqref{dy:f16} are consequently a
finite, exact specification of one integer-coefficient polynomial.

\begin{theorem}[Ten-positive equation]
\label{dy:main}
With the fixed coefficient binding and source calibration of
Section~\ref{sec:calibration}, the polynomial $N$ has a zero in
$\Pos^{10}$ if and only if ZFC has a finite syntactic refutation.
The reduction is ordinary mathematics formalizable in ZFC.  Thus the
source calibration gives
\[
 \mathrm{ZFC}\vdash
 \bigl[(\forall\mathbf x\in\Pos^{10})\,N(\mathbf x)\ne0\bigr]
       \longleftrightarrow\Con(\mathrm{ZFC}).
\]
The exact total degree of $N$ is $706$.  The supplied literal circuit has
$133$ binary operations: $79$ multiplications, $33$ additions and $21$
subtractions, including operations with small numerals.
\end{theorem}

\subsection{The norm identity}
\label{dy:normproof}

We include the conjugate-product calculation because it specifies exactly
what the single polynomial equation asserts.

\begin{lemma}[Sixteen conjugates]
\label{dy:normlemma}
In the polynomial ring over $\Z$,
\begin{equation}
 \mathcal F_{16}(x^2;r_1^2,r_2^2,r_3^2,r_4^2)
 =\prod_{\varepsilon\in\{1,-1\}^4}
       (x-\varepsilon_1r_1-\varepsilon_2r_2
          -\varepsilon_3r_3-\varepsilon_4r_4).
 \label{dy:normidentity}
\end{equation}
Moreover, $\mathcal F_{16}$ is homogeneous of degree $8$ in its five
arguments and has leading coefficient $1$ in $s$.
\end{lemma}

\begin{proof}
Put $a=r_1^2$, $b=r_2^2$, $c=r_3^2$ and $d=r_4^2$.
Pairing the signs of $r_1$ and $r_2$, then the sign of $r_3$, gives the
three-root identity
\[
 \prod_{\varepsilon_1,\varepsilon_2,\varepsilon_3}
   (y-\varepsilon_1r_1-\varepsilon_2r_2-\varepsilon_3r_3)
 =\bigl[(y^2-\sigma_1)^2-4\sigma_2\bigr]^2
       -64\sigma_3y^2.
\]
For example the first pairing is
$(y^2-a-b)^2-4ab$; substituting $y-r_3$ and $y+r_3$ and multiplying yields
the displayed expression.

Now fix the sign $\varepsilon$ of $r_4$ and set
$y=x-\varepsilon r_4$, $s=x^2$ and $r=2xr_4$.  Then
$y^2=U-\varepsilon r$ and $r^2=4sd$.  Consequently
\[
 (y^2-\sigma_1)^2-4\sigma_2=E-2\varepsilon rP.
\]
The three-root product becomes
\[
 (E-2\varepsilon rP)^2-64\sigma_3(U-\varepsilon r)
       =G_*-4\varepsilon rO_*.
\]
Multiplication over the two values of $\varepsilon$ gives
$G_*^2-16r^2O_*^2=G_*^2-64sdO_*^2$, as claimed.  The product identity also
shows the homogeneity and the coefficient of $s^8$.
\end{proof}

For later use define positive real roots, whenever the positive input is
fixed, by
\begin{equation}
 \tau=\sqrt{D_\tau},\qquad f=\sqrt{D_f},\qquad
 y=\sqrt{D_y/\delta},\qquad D=\sqrt{D_m},\qquad d=D/H.
 \label{dy:roots}
\end{equation}
Their existence for \emph{every} positive input is checked below, before
any equation or transport sign is assumed.  Lemma~\ref{dy:normlemma}
then yields
\begin{equation}
 N=\delta^8M^{16}
    \prod_{\varepsilon\in\{1,-1\}^4}
    \left(\frac ZM-4\varepsilon_1\tau-2\varepsilon_2f
                       -\varepsilon_3y-8\varepsilon_4d\right).
 \label{dy:factorization}
\end{equation}
Direct division of the definition of $Z$ gives the explanatory identity
\begin{align}
 \frac ZM={}&n+V+2C_f+4C_\tau+\frac{8(X+a_-c)}H
       +\frac{Y}{q^3}-\frac{WB_t}{q-1}-\frac e{\mu^2}.
 \label{dy:carrieridentity}
\end{align}
Only the polynomial definition~\eqref{dy:carrier}, not division, is used
in evaluating $N$.

\subsection{Square classes separated by powers of two}
\label{dy:separation}

\begin{lemma}[Dyadic separation]
\label{dy:dyadic}
Let $r_1,\ldots,r_m$ be positive real numbers with rational squares and
$v_2(r_i^2)=0$.  Let $h_1,\ldots,h_m$ be distinct integers.  If, for some
signs $\epsilon_i$, the sum
$\sum_i\epsilon_i2^{h_i}r_i$ is rational, then all the $r_i$ are rational.
\end{lemma}

\begin{proof}
The extension $K=\Q(r_1,\ldots,r_m)$ is a finite multiquadratic Galois
extension.  Its group sends each $r_i$ to $\pm r_i$, defining a character
$\chi_i$.  Two roots have the same character exactly when their ratio is
rational; a character is trivial exactly when its root is rational.
If some nontrivial character $\chi$ occurs, apply its projection
\[
  \frac1{|\operatorname{Gal}(K/\Q)|}
  \sum_{g\in\operatorname{Gal}(K/\Q)}\chi(g)g
\]
to the rational relation.  It yields
\[
 \sum_{\chi_i=\chi}\epsilon_i2^{h_i}r_i=0.
\]
Divide by any root in this nonempty group.  The ratios are rational and
have $2$-adic valuation zero, because their squares have valuation zero.
The resulting rational summands have pairwise distinct valuations $h_i$.
The term of least valuation cannot cancel.  Equivalently, clear the odd
denominators and divide by the least power of two; one summand is odd and
all the others are even.  This contradiction excludes every nontrivial
character.
\end{proof}

The lemma uses valuations of rational ratios only.  It does not require
an extension of the rational $2$-adic valuation to arbitrary irrational
real numbers, and it imposes no relative-size assumption on the roots.

\begin{lemma}[All-input positivity and units]
\label{dy:allinput}
For every input in $\Pos^{10}$, $\delta,D_y,H$ are positive and the four
radicands in~\eqref{dy:roots} are positive rational numbers of $2$-adic
valuation zero.
\end{lemma}

\begin{proof}
The number $\beta$ is odd, as is $J$, so $q$ is even and $q\ge16$.
Thus $\mu$ is odd, $\lambda$ is even, $A$ is odd, $\Delta$ is even,
$V_0$ is even and $H$ is odd.  Also $k\ge2$, $c\ge3$, $I\ge c^2$ and
$\Delta>\mu^2>1$.  In particular $Q>1$ and $\delta>0$.

If the integer $b=cT'-R$ is zero, then $V=-c$ and $|V|\ge3$.
If $b\ne0$, then
\[
 |V|=|D_fb-c|\ge D_f-c>1,
\]
since $D_f\ge1+\Delta c^4$.  Hence $D_y=QV^2-1>0$ in both cases.
Every numerator and denominator in
\[
 D_\tau,\qquad D_f,\qquad \frac{D_y}{\delta},\qquad
 \frac{D_m}{H^2}
\]
is odd: $D_\tau=1+V_0(V_0+1)k^2$, $D_f=1+\Delta I^2$,
$Q=\Delta^2I^2$ is even, and $D_m=1+\Delta c^2$.
All four rational numbers are therefore units at $2$.
\end{proof}

\begin{lemma}[Adjacent-coefficient descent]
\label{dy:descent}
Let $Q>1$ and $V$ be integers.  If $y\in\Q$ satisfies
\[
 QV^2-(Q-1)y^2=1,
\]
then $y$ is an integer.
\end{lemma}

\begin{proof}
Replace $V,y$ by their absolute values.  Neither can be zero.
Write $y=u/m$ in lowest terms, with $m>0$.  The equation implies
$m^2\mid Q-1$.  When $V=1$, it gives $y=1$.
For $V>1$ define
\[
 V'=(2Q-1)V-2(Q-1)y,\qquad
 y'=(2Q-1)y-2QV.
\]
The first quantity is an integer because $m\mid Q-1$.  The identity
$(2Q-1)^2-4Q(Q-1)=1$ gives
\[
 Q(V')^2-(Q-1)(y')^2=1.
\]
Furthermore
\[
 (Q-1)(y^2-V^2)=V^2-1>0
\]
shows $y>V$, whereas
\[
 (QV)^2-((Q-1)y)^2=QV^2+Q-1>0
\]
shows $QV>(Q-1)y$.  These inequalities imply $-V<V'<V$.
The new equation excludes $V'=0$.  Since $2Q-1\equiv1\pmod m$,
the reduced denominator of $y'$ is still $m$.
Induction on the positive integer $|V|$, using $|V'|<V$, reaches $V=1$.
Thus $m=1$.
\end{proof}

If $N=0$, equation~\eqref{dy:factorization} gives a rational signed
sum of the roots with weights $(1,2,4,8)$ in the order $(y,f,\tau,d)$.
Lemmas~\ref{dy:allinput} and~\ref{dy:dyadic} make all four roots rational.
The roots $\tau,f,D$ have integer squares and hence are integers.
Lemma~\ref{dy:descent} applied to
$QV^2-(Q-1)y^2=1$ makes $y$ an integer as well.  This step is valid even
when $V$ is still negative: the descent can use $|V|$.

\subsection{Soundness: extracting transport and the canonical data}
\label{dy:soundness}

Assume $N=0$ at a positive input, and choose signs from a vanishing factor
of~\eqref{dy:factorization}.  The first congruence is obtained before
transport positivity.  Since $q$ is even,
\[
 \gcd(\mu^2,q^3(q-1))=1,
 \qquad H=4\lambda(a+2)-1\equiv-1\pmod\mu.
\]
Thus $\mu^2$ is coprime to $q^3(q-1)H$.  Integer-valuedness of the
signed roots in~\eqref{dy:carrieridentity} gives
\[
 \frac{8(X+a_-c-\varepsilon_4D)}H+\frac Y{q^3}
           -\frac{WB_t}{q-1}-\frac e{\mu^2}\in\Z.
\]
Clearing the other denominators therefore gives $\mu^2\mid e$.
Put
\begin{equation}
 \rho=\frac e{\mu^2}\in\Z,\qquad
 g=\frac{\varepsilon_4D-X-a_-c}{H},\qquad
 s_0=\frac Y{q^3}>0,\qquad t=\frac{B_t}{q-1}.
 \label{dy:fractionnames}
\end{equation}
At this stage $g,s_0,t$ need not be integers.  The carrier identity is
\begin{align}
 n={}&4(\varepsilon_1\tau-C_\tau)
       +2(\varepsilon_2 f-C_f)+(\varepsilon_3y-V)
       +8g-s_0+Wt+\rho.
 \label{dy:nequation}
\end{align}

Use the standard Pell notation
\[
 \chi_B(r)+\psi_B(r)\sqrt{B^2-1}
       =(B+\sqrt{B^2-1})^r.
\]
The integer roots give, for positive indices,
\begin{align}
 \tau&=\chi_P(n_0),&k&=2\psi_P(n_0),&P&=2V_0+1,\notag\\
 D&=\chi_A(p),&c&=\psi_A(p),\notag\\
 f&=\chi_A(m),&I&=\psi_A(m).
 \label{dy:pellindices}
\end{align}
For the first equation the fundamental solution of
$x^2-V_0(V_0+1)k^2=1$ is $(2V_0+1,2)$: $k=1$ would require a square
strictly between $V_0^2$ and $(V_0+1)^2$.
If $n_0=1$, then $k=2$, $\eta=\zeta=1$ and $c=2Y+1$.
But $\psi_A(1)=1$ and $\psi_A(p)\ge2A>2Y+1$ for $p\ge2$.
Consequently $n_0\ge2$, and
\begin{equation}
 c>k\ge4P>8X\lambda^2Y^2>8q^5.
 \label{dy:earlyc}
\end{equation}

The mask bound~\eqref{dy:ports} gives
\[
 0<(M_C+qM_F)J<2q^2.
\]
If $F\ge q$, then $G\le\alpha$; if $F<q$, then
$G<(q-1)^2+\alpha$.  In either case
\begin{equation}
 R<q^4+(\alpha+2)q^2.
 \label{dy:Ronesided}
\end{equation}
Together with~\eqref{dy:earlyc}, this proves $c(T+\alpha)>R$.
Hence $b\ge1$, $V=D_fb-c>1$ and $y>V$.
No conclusion about $B_t$ has yet been used.

\begin{lemma}[Quadratic transport bound]
\label{dy:transport}
Every positive zero of $N$ satisfies $B_t\ge0$.
\end{lemma}

\begin{proof}
Before any sign restriction on $R$, the definition of $R$ gives
\begin{equation}
 |R|<q^4+Fq^3+(\alpha+2)q^2.
 \label{dy:Rabsolute}
\end{equation}
As $q\ge16$ and $\mu^2>q^4$, it follows that
\begin{equation}
 \frac{|R|}{\mu^2}<F+\alpha+1,\qquad
 |\rho|\le\frac{k+1+|R|}{\mu^2}<k+F+\alpha+1.
 \label{dy:rhoestimate}
\end{equation}
The root centers satisfy the one-sided estimates
\begin{equation}
 \tau-C_\tau<k,\qquad f-C_f<0,\qquad 8g<k.
 \label{dy:rootupper}
\end{equation}
For the first inequality, square the comparison with $(V_0+1)k$.
For the second, $D_f=A^2I^2-I^2+1<A^2I^2$, since $I>1$.
For the third, $D<Ac$ implies $g<2c/H$; now $c<2kY$ and $H>32Y$.
These upper bounds apply to the corresponding signed roots in
\eqref{dy:nequation} as well.

Put $u=y-V$.  Since $y>V>1$,
\begin{align*}
 u&=\frac{V^2-1}{(Q-1)(y+V)}
       <\frac{V}{2(Q-1)}
       <\frac{D_fb}{2(Q-1)}<\frac b\Delta.
\end{align*}
The last inequality is equivalent to
$\Delta D_f=Q+\Delta<2Q-2$, which holds because
$Q=\Delta^2I^2>\Delta+2$.  Also $\Delta>\mu^2$ and
$c<2kY<k\Delta$.  Thus
\begin{equation}
 u<kT'+F+\alpha+1.
 \label{dy:gapupper}
\end{equation}

If the integer $B_t$ were negative, it would be at most $-1$, and
$Wt\le-W/(q-1)<-L$.  Substituting
\eqref{dy:rhoestimate}--\eqref{dy:gapupper} into
\eqref{dy:nequation}, and using
$\varepsilon_3y-V\le y-V$ and $s_0>0$, yields
\[
 n<kT'+6k+2F+2\alpha+2-L<0.
\]
Indeed
\begin{align*}
 L-(kT'+6k+2F+2\alpha+2)
  &=(F+1)T'+4F+4-2\alpha\\
  &=(F+1)T+(F-1)\alpha+4F+4>0.
\end{align*}
This contradicts the stipulated $n\ge1$.
\end{proof}

Now $(K_0+w)C\ge F>0$, so $C\ge1$.  The source outer estimates give
\begin{equation}
 3q<R<q^4-2.
 \label{dy:Router}
\end{equation}
For completeness, write
$G=(q-1)C+q\alpha$ and use $C+\alpha=q-F\le q-1$.
Its lower bound and the positive mask term imply the left inequality;
$G\le q^2-q$ and the mask bound below $2q^2$ imply the right one.

We record how the inherited index arguments apply; this is also a check
that deleting the factor $\mu$ from $I$ changes no hypothesis.
The first Pell congruence is
\[
 k\equiv2n_0\pmod{4V_0},\qquad \mu^2\mid4V_0.
\]
Together with $\mu^2\mid k-R-1$ and~\eqref{dy:Router}, it gives
$2n_0\ge R+1$.  Equations~\eqref{dy:earlyc} and~\eqref{dy:Router}
give $c>2R$; the main Pell equation gives $c>2p$.
The rank lemma requires $c^2\mid I$, which holds exactly as written in
\eqref{dy:auxiliary}.  It gives $c\mid m$ and $m>2p$.
With $S=\Delta I$ the auxiliary data satisfy
\[
 S^2=\Delta(f^2-1)=Q,\qquad c\mid S,\qquad
 V\equiv-c\pmod f,\qquad V\equiv-R\pmod c.
\]
The source index lemma therefore gives $p=R$.
Since $P>A$, monotonicity of Pell ordinates gives $n_0<R$.
The congruence $2n_0\equiv R+1\pmod{\mu^2}$, together with
$R<q^4<\mu^2$, forces $2n_0=R+1$.  The fixed-minus parity lemma then
implies $R\equiv3\pmod4$.

Set $r=(R-1)/2$.  The parent ratio lemma is valid for every integer
scaling $\lambda\ge1$ and hence for $\lambda=2\mu$.  In its notation,
\begin{equation}
 \xi=\frac{(X+1)^{2r}}{X^r},\qquad
 \xi<\frac{\psi_A(2r+1)}{\psi_P(r+1)}
       <\xi\left(1+\frac3a\right)^{2r}.
 \label{dy:ratio}
\end{equation}
The defining interval $kY<c<k(Y+1)$ gives
\begin{equation}
 Y>\frac{X^r}{3},\qquad a>\frac{X^{r+1}}3,\qquad r\ge24.
 \label{dy:Ybound}
\end{equation}
Before the exponent has been determined, clearing the three remaining
fractions gives
\[
 H\mid8q^3(q-1)(X+a_-c-\varepsilon_4D).
\]
Because $H$ is odd, the factor $8$ cancels.  The main projection identities
are
\[
 D-a_-c\equiv2^R\pmod H,\qquad
 D+a_-c\equiv2^{-R}\pmod H.
\]
The parent Archimedean estimate, using~\eqref{dy:Ybound}, is
\begin{equation}
 q^4,2^{R+1}X<X^{r+1}<3a<H.
 \label{dy:archimedean}
\end{equation}
If $\varepsilon_4=-1$, multiplication by $2^R$ would make
$q^3(q-1)(2^RX+1)$ a positive multiple of $H$ smaller than $H$.
This is impossible.  If $\varepsilon_4=+1$, the same estimate bounds
$|q^3(q-1)(X-2^R)|$ strictly below $H$, so divisibility gives
\begin{equation}
 \varepsilon_4=+1,\qquad X=2^R.
 \label{dy:exponent}
\end{equation}
In particular $g$ is an integer.  The remaining partial fractions now give
$q^3\mid Y$ and $q-1\mid B_t$.

The ratio estimate~\eqref{dy:ratio} and the half-binomial lemma yield
\begin{equation}
 Y=\frac12\sum_{h=0}^r\binom{2r}{r+h}X^h.
 \label{dy:halfbinomial}
\end{equation}
To see the uniqueness mechanism, write
\[
 \xi=M_0+\theta,\qquad
 M_0=\sum_{h=0}^r\binom{2r}{r+h}X^h,\qquad0<\theta<\tfrac14.
\]
Here $M_0$ is even because $X=2^{2r+1}$ and the central binomial
coefficient is even.  The quantitative ratio estimate puts $c/k$ in
$(M_0/2,M_0/2+1)$, while its defining interval is $(Y,Y+1)$.
These two integer parts must agree.  The explicit estimates and source
reconstruction are the ones in~\cite[Sections 6.17--6.18]{parent816};
only their arbitrary positive scaling parameter has changed.

Finally keep $J,F,\alpha,w,Y$ and set
$s=Y/q^3$, $z=B_t/(q-1)+1$.
Reconstruct fresh \emph{unscaled} Pell witnesses with main parameter
$a+2$ and first parameter $2XY^2+1$, using
\eqref{dy:exponent}, $R\equiv3\pmod4$ and~\eqref{dy:halfbinomial}.
This is the source reconstruction of~\cite[Section 6.18]{parent816}:
it requires only the recovered canonical data and supplies all fifteen
positive witnesses of the calibrated source system.
Theorem~\texttt{thm:source} of~\cite{source61} gives a finite syntactic
refutation of ZFC.  This proves soundness relative to the calibrated
interface.

\subsection{Completeness and the positive sum coordinate}
\label{dy:completeness}

Suppose ZFC has a finite syntactic refutation.  The calibrated source
provides a canonical positive solution with
\[
 q=B^N\text{ for odd }N>1,\quad R\equiv3\pmod4,\quad
 X=2^R,\quad C>0,
\]
and $Y$ equal to~\eqref{dy:halfbinomial}.  Its odd $J$ gives the positive
integer $j=(J-1)/2$.
The parent ratio argument, valid for arbitrary positive integer scaling,
constructs main and first Pell values for $\lambda=2\mu$:
\[
 c=\psi_A(R),\quad D=\chi_A(R),\quad
 k=2\psi_P((R+1)/2),\quad \tau=\chi_P((R+1)/2).
\]
Set $\eta=c-kY$, $\zeta=k-\eta$; both are positive.
The quantities
\[
 \rho=(k-R-1)/\mu^2,\qquad
 g=(D-X-a_-c)/H,\qquad t=B_t/(q-1)
\]
are integers with $\rho,g\ge1$ and $t\ge0$, as in the parent
completeness proof.

For the strong Pell pair use the smaller index
\begin{equation}
 m=Rc,\qquad I=\psi_A(m),\qquad f=\chi_A(m).
 \label{dy:strongcomplete}
\end{equation}
The multiplication formula and the standard Pell congruence give
\[
 \psi_A(Rc)=c\psi_D(c),\qquad
 \psi_D(c)\equiv cD^{c-1}\pmod{D^2-1}.
\]
Since $c^2\mid D^2-1$, this implies $c\mid\psi_D(c)$ and therefore
$c^2\mid I$.  Thus $i=I/c^2$ is a positive integer.
No divisibility by $\mu$ is needed.

Let $S=\Delta I$ and $Q=S^2$.  The explicit solution in the auxiliary
completeness lemma has
\begin{align}
 V^*&=\frac{\chi_S(R)}S=Q_v(S^2),&y&=\psi_S(R),&v&=(R-1)/2,\notag\\
 T^*&=\frac{V^*+c+Rf^2}{cf^2},&
 Q(V^*)^2-(Q-1)y^2&=1.
 \label{dy:auxcomplete}
\end{align}
The inherited lemma proves that $T^*$ is an integer.  The associated odd
Pell polynomial obeys
\[
 Q_v(z)\ge(4z-3)^v\quad(z\ge1).
\]
Since $v\ge2$ here,
\begin{equation}
 V^*\ge(4Q-3)^v\ge(3Q)^v\ge9Q^2.
 \label{dy:Vlarge}
\end{equation}
As $c<S$ and $f^2=D_f<Q$, it follows that
\[
 T^*>\frac{V^*}{S^3}\ge9S>\alpha.
\]
The last comparison uses the canonical $C>0$, hence
$\alpha<q<S$.  Define $T=T^*-\alpha\ge1$; then $T'=T^*$ and
$V=V^*$ as required.

There is one new completeness issue: the center $C_f=AI$ lies above
$f$, so its contribution to $n$ is negative.  The auxiliary gap has ample
room to absorb it, with an explicit bound.

\begin{lemma}[Canonical auxiliary gap]
\label{dy:gaplower}
For the witnesses just constructed, $u=y-V$ satisfies
\[
 u>2I+Y+1.
\]
\end{lemma}

\begin{proof}
Since $V\ge2$, $Q\ge3$ and $V<y<2V$,
\[
 u=\frac{V^2-1}{(Q-1)(y+V)}
      >\frac V{4Q}
      \ge\frac{9S^2}{4}>3S+1.
\]
For the first strict inequality one may use
$V^2-1\ge\tfrac34V^2$ and
$(Q-1)(y+V)<3QV$.  The second is~\eqref{dy:Vlarge}.
Finally $Y<c<I<S$, so $3S+1>2I+Y+1$.
\end{proof}

Define the integer
\begin{equation}
 n=Wt+\rho+8g+4(\tau-C_\tau)+2(f-C_f)+u-Y/q^3.
 \label{dy:ncomplete}
\end{equation}
The terms $Wt,\rho,g$ are nonnegative or positive, and $\tau>C_\tau$.
Moreover $f>(A-1)I$, so $f-C_f>-I$, and $Y/q^3\le Y$.
Lemma~\ref{dy:gaplower} therefore gives $n>0$, hence $n\in\Pos$.
Substitution into~\eqref{dy:carrieridentity} yields
\[
 Z/M=4\tau+2f+y+8d.
\]
The all-positive sign factor in~\eqref{dy:factorization} vanishes, proving
$N=0$.  This establishes completeness and the semantic part of
Theorem~\ref{dy:main}.

\subsection{A symbolic proof of the exact degree}
\label{dy:degree}

The literal circuit gives the conservative syntactic degree $832$.
That estimate ignores cancellations caused by the three centers.  We
prove the exact total degree using a formal parameter, so no numerical
specialization is needed for the upper bound.

Replace every unknown $x$ by $tx$, keeping the direction variables
algebraically independent, and work over their rational function field
$K$.  Degree below means degree in $t$ at infinity.  In $K((t^{-1}))$
each relevant square root has a branch whose leading term is its center:
this follows by the binomial expansion of $(1+h)^{1/2}$ when $h$ has
negative degree.  For $y$, choose the branch with leading term $V$.
These formal branches need not be positive real branches; the product
over all signs is independent of that choice.

The register and error degrees are
\begin{center}
\begin{tabular}{@{}lr@{}}
\toprule
Quantity & Degree\\
\midrule
$\mu,\lambda$ & $2$\\
$A,H,\Delta$ & $5,5,10$\\
$I,Q,\delta$ & $5,30,30$\\
$D_f,V,M$ & $20,24,13$\\
$C_f,C_\tau$ & $10,9$\\
$\tau-C_\tau$ & $1$\\
$f-C_f$ & $0$\\
$y-V$ & $-6$\\
$d$ & $2$\\
$L,W,B_t$ & $2,3,2$\\
$Z/M-V-2C_f-4C_\tau$ & $4$\\
\bottomrule
\end{tabular}
\end{center}
For example,
\[
 f^2-C_f^2=1-I^2
\]
has degree $10$, while $f+C_f$ has degree $10$.
Similarly
\[
 (y-V)(y+V)=\frac{V^2-1}{Q-1}
\]
gives degree $48-30-24=-6$ for $y-V$.
The degree-$4$ leading term of the remaining carrier is $-L_2(B_t)_2$,
coming from $-WB_t/(q-1)$; its other terms have degree at most $2$.

The sixteen factors in~\eqref{dy:factorization} have the following
nonzero leading terms:
\begin{center}
\begin{tabular}{@{}p{0.47\linewidth}rrl@{}}
\toprule
Signs & Count & Degree & Leading term\\
\midrule
$\varepsilon_3=-1$ & $8$ & $24$ & $2V_{24}$\\
$\varepsilon_3=1$, $\varepsilon_2=-1$ & $4$ & $10$ & $4(C_f)_{10}$\\
$\varepsilon_3=\varepsilon_2=1$, $\varepsilon_1=-1$
    & $2$ & $9$ & $8(C_\tau)_9$\\
$\varepsilon_3=\varepsilon_2=\varepsilon_1=1$
    & $2$ & $4$ & $-L_2(B_t)_2$\\
\bottomrule
\end{tabular}
\end{center}
Subscripts here denote leading homogeneous forms.  Their product has
exact degree
\[
 8\cdot24+4\cdot10+2\cdot9+2\cdot4=258.
\]
The prefactor $\delta^8M^{16}$ has degree
$8\cdot30+16\cdot13=448$.  Thus
\begin{equation}
 \deg N=258+448=706.
 \label{dy:exactdegree}
\end{equation}

For an explicit check that the leading form is nonzero, put
\begin{align*}
 q_0&=2\beta j,& k_0&=\eta+\zeta,& c_0&=k_0Y,\\
 \lambda_0&=2q_0^2,& A_0&=\lambda_0wq_0Y,&\Delta_0&=A_0^2,\\
 I_0&=ic_0^2,& Q_0&=\Delta_0^2I_0^2,\\
 V_*&=-(\Delta_0I_0^2)q_0^3(q_0-F),\\
 C_{f,0}&=A_0I_0,&
 C_{\tau,0}&=wq_0\lambda_0^2Y^2k_0,\\
 M_0&=4q_0^8A_0,& L_0&=(k_0+F)(T+\alpha),&
 B_{t,0}&=w(q_0-F-\alpha).
\end{align*}
The degree-$706$ component is
\begin{equation}
 2^{22}Q_0^8M_0^{16}V_*^8C_{f,0}^4C_{\tau,0}^2
                    (L_0B_{t,0})^2.
 \label{dy:leadingform}
\end{equation}
Every factor is a nonzero polynomial.  More concretely, on the diagonal
where all ten unknowns equal $t$, put $q_0=2\beta$ as a fixed integer.
The coefficient of $t^{706}$ is
\begin{equation}
 9\cdot2^{202}q_0^{366}(q_0-1)^8(q_0-2)^2>0.
 \label{dy:diagonal}
\end{equation}
This explicit coefficient supplies a lower-degree witness for the actual
fixed coefficient tuple.  The three other coefficient ports occur only
in lower-degree terms.

\subsection{Circuit form, verification and limitations}
\label{dy:verification}

The proof uses $W$ as a convenient abbreviation.  The literal circuit
does not materialize it: it distributes
\begin{equation}
 \frac{WB_t}{q-1}=LB_t+\frac{B_t}{q-1}.
 \label{dy:weightsave}
\end{equation}
Thus the circuit subtracts $LB_t$ directly from
$n+V+2C_f+4C_\tau$, and uses $B_t$ alone in the residual fractional
numerator.  This exact identity saves one binary operation and leaves the
polynomial unchanged.  The final count is $79+33+21=133$; all rows are
live and every unknown and coefficient is used.

The supplied build program checks the four compressed and raw coefficient
hashes, their bit lengths and bounds, the circuit topology, liveness and
operation counts, and deterministic regeneration of the JSON circuit.
The audit evaluates the displayed formulas independently at modular
points and computes complete degree-bounded univariate restrictions
modulo two primes.  These restrictions have degree $706$ and the leading
coefficient~\eqref{dy:leadingform}; the diagonal restrictions also match
\eqref{dy:diagonal}.  Additional exact finite cases check the parity and
majorant premises and the lower bound on the canonical auxiliary gap.
These are implementation checks supporting the proof, rather than a
substitute for it.

The reduction is a measurable improvement in degree over the pinned
parent, with a modest increase in circuit length.  It does not claim the
least possible degree, the fewest possible unknowns, or the shortest
possible equation.  It retains the inherited K215/native compiler route
and its stated formalization boundary.  For a formulation over
nonnegative integers, replacing each unknown $x$ by $x+1$ preserves the
exact degree and adds at most ten additions to the literal evaluator.

\section{A direct quartic in 780 unrestricted integers}\label{dc:section}

The second construction reduces the direct $C_{222}$ quartic of
\cite{direct407} from 888 to 780 integer unknowns.  It keeps degree four,
reduces the collected support from 10,156 to 9,628 monomials, and leaves
the maximum absolute expanded coefficient unchanged at
$2\,445\,790\,974\,704$.  Its new ingredient is an order of proof:
the signs of the private power quotients are recovered \emph{after}
the power graphs have been shown correct.  This removes independent
sign guards without weakening the final arithmetic specification.

\subsection{Exact object and inherited source}

Let $r_0,\ldots,r_{433}$ be the literal quadratic rows in this component's
\texttt{equations/equation.json}, using its listed order of 780 variables.
Define the fixed integer polynomial
\begin{equation}\label{dc:polynomial}
 P_{780}(x_0,\ldots,x_{779})=\sum_{j=0}^{433}r_j(x)^2.
\end{equation}
All 780 variables range over $\Z$.  The complete collected polynomial is
also supplied as \texttt{expanded\_polynomial.json}; this is an ordinary
polynomial with literal integer coefficients and fixed exponents.
The coefficient of $x_0^4$ is 3, so its degree is exactly four.

The pinned immediate source is the archive
\texttt{ConZFC\_407\_888\_Quartics\_2026-10-24.zip}
at ProveIt commit
\texttt{ec80841c07e91be8956ce6ed4af9a9c0b782d3d9}.
It is retained byte for byte, with its original notices; the separate
\texttt{SOURCE\_PINS.json} records the archive digest and the exact members
used.  Write $S(y)=0$ for its 606 quadratic source rows on $y\in\N^{690}$.
The inherited arithmetic and logical calibration is
\begin{equation}\label{dc:source}
 \bigl(\exists y\in\N^{690}: S(y)=0\bigr)
 \ \longleftrightarrow\ C_{222}(1)
 \ \longleftrightarrow\
 \exists p\;\operatorname{Prf}_{\mathrm{ZFC}}(p,\bot).
\end{equation}
The source report identifies the inherited calibration theorems
\begin{quote}\small\ttfamily
KEndpoint.c222\_iff\_zf\_refutation\\
ConstructibleChoice.zf\_refutation\_iff\_zfc\_refutation
\end{quote}
Their calibration snapshot is
\begin{center}\small
\texttt{08200eb10fde2db408c69dffd840d6920a569b1f}.
\end{center}
Here we prove the new existence equivalence with $S$; the proof of
\eqref{dc:source} is inherited from that explicitly identified source.

\subsection{Guarded signs and a staged inference certificate}

For a selected source coordinate $y_v$, three private integer unknowns
enforce the quadratic guard
\begin{equation}\label{dc:guard}
 y_v=a_v^2+b_v^2+c_v(c_v+1).
\end{equation}
This forces $y_v\ge0$ for all integers $a_v,b_v,c_v$.  Conversely, every
nonnegative integer has such a representation: apply Legendre's
three-square theorem to $4y_v+1$, whose residue modulo four forces one
of the three squares to be odd and the other two even, then subtract
one and divide by four.  The parent used 167 such guards.  We use 131.

Ordinary sign propagation is performed on the unchanged source rows.
An oriented row is factored exactly as
$a(\widehat y)y_v-b(\widehat y)=0$, with neither $a$ nor $b$ involving
$y_v$.  Previously proved lower bounds $a\ge1$ and $b\ge0$ imply
$y_v\ge0$; if $b\ge1$, integrality implies $y_v\ge1$.
The checker establishes these bounds from sums of nonnegative
monomials and exact affine squares.  Odd powers need an established
nonnegative factor; even powers need no sign premise.  A monomial
whose factors are all at least one contributes its coefficient to
the constant lower bound.  This last rule can certify a polynomial
with a negative constant.  Every premise must have been established
earlier in the certificate.

Starting from the 131 guards, these elementary inferences prove 623
of the 690 source coordinates nonnegative.  They also prove that every
power base and every exponent is at least one, and prove all sign
premises for the frame below.  In particular, a factorial output used
as a power base receives its strict sign from literal arithmetic rows;
the proof does not first appeal to factorial correctness.

\subsection{The frame bounds do not require power-quotient signs}

There are 23 power calls $(n,k,m)$, nine distinct bases, and 47 distinct
raw atoms occurring as a base, exponent, or proposed output.  Denote
that set of atoms by $\mathcal A$.  The exact shared frame is
\begin{equation}\label{dc:frame}
 W=2+\sum_{a\in\mathcal A}a,
 \qquad Z=z+1,
 \qquad A^2-W^3(W+2)Z^2=1.
\end{equation}
At this point all atoms in $\mathcal A$ are known nonnegative,
$n,k\ge1$, $A\ge0$ and $Z\ge1$.  Consequently
$W\ge2$, $1\le n,k\le W$, $0\le m\le W$, and $A>W^2$.

For $a\ge2$, define the Pell coordinates by
\begin{align}
 X_0(a)&=1,&Y_0(a)&=0,\nonumber\\
 X_{j+1}(a)&=aX_j(a)+(a^2-1)Y_j(a),&
 Y_{j+1}(a)&=X_j(a)+aY_j(a).\label{dc:pell}
\end{align}
The norm in \eqref{dc:frame} is also
$A^2-((W+1)^2-1)(WZ)^2=1$.
Pell-solution classification therefore gives
$A=X_j(W+1)$ and $WZ=Y_j(W+1)$ for some $j>0$.
Reducing \eqref{dc:pell} modulo $W$ gives
$X_j(W+1)\equiv1$ and $Y_j(W+1)\equiv j\pmod W$.
Thus $W\mid j$ and $j\ge W$, whence
\begin{equation}\label{dc:exponentialbound}
 A\ge(W+1)^j\ge(W+1)^W>W^W\ge n^k.
\end{equation}

For each distinct base the literal system defines
\begin{equation}\label{dc:modulus}
 T_n=2An-n^2-1,
 \qquad D=A^2-1.
\end{equation}
Their signs have not been independently guarded.  Since
$1\le n\le W<A$, the function $2An-n^2-1$ is increasing on this
integer interval.  Therefore
\begin{equation}\label{dc:modulusbound}
 T_n\ge2A-2>W^W\ge n^k,
 \qquad T_n>W\ge m,
 \qquad D\ge1.
\end{equation}
We have recovered the nine modulus signs and the discriminant sign
using only the frame and earlier elementary facts.  Elementary
propagation from these facts raises the number of nonnegative source
coordinates to 644.  No private power quotient was used in this step.

\subsection{Index certification precedes power recovery}

Fifteen seed exponents have individual witnesses
$X_i,Y_i,S_i,R_i,e_i,s_i,r_i$, with a shared $U,V,B$ and auxiliary
quotients $b,v,q$.  The exact equations are
\begin{align}
 Y_i&=k_i+e_i,& X_i^2-(A^2-1)Y_i^2&=1,\nonumber\\
 S_i^2-(B^2-1)R_i^2&=1,& S_i&=X_i+Us_i,\nonumber\\
 R_i&=k_i+4Y_i r_i,& U^2-(A^2-1)V^2&=1,\label{dc:index}\\
 B&=A+Ub,&
 V&=\left(\prod_{i=1}^{15}Y_i^2\right)(v+1),\nonumber\\
 B&=1+4V(q+1).&&\nonumber
\end{align}
Every displayed witness has already been proved nonnegative, and
$k_i\ge1$.  Hence $Y_i\ge1$, $V\ge1$, $B\ge5$,
$Y_i^2\mid V$, and $4Y_i\mid(B-1)$.
For each $i$ these equations satisfy the hypotheses of the inherited
Matiyasevich Pell-index criterion
\texttt{PellRepresentation.criterion\_iff\_pell}.  Thus
\begin{equation}\label{dc:indexidentified}
 X_i=X_{k_i}(A),\qquad Y_i=Y_{k_i}(A).
\end{equation}
The three remaining requested pairs are derived by shifts $s=1$ or
$s=2$, using the exact polynomial identities
\begin{equation}\label{dc:addition}
 X_{k+s}=X_kX_s+DY_kY_s,
 \qquad Y_{k+s}=X_kY_s+Y_kX_s,
\end{equation}
where $(X_1,Y_1)=(A,1)$ and $(X_2,Y_2)=(2D+1,2A)$.
Their gate identities and exponent shifts are checked literally.
All required Pell pairs are therefore identified before using the
sign of any power quotient.

\subsection{The quotient sign follows from the exact power}

\begin{lemma}[Delayed power-quotient sign]\label{dc:quotientlemma}
Suppose $A\ge2$, $n,k\ge1$, $T=2An-n^2-1>0$,
$0\le m<T$ and $n^k<T$.  Suppose further that
$X=X_k(A)$, $Y=Y_k(A)$, and that an integer $h$ satisfies
\begin{equation}\label{dc:power}
 X+nY=AY+m+Th.
\end{equation}
Then $m=n^k$ and
\begin{equation}\label{dc:hformula}
 h=\sum_{j=0}^{k-1}n^{k-1-j}Y_j(A)\ge0.
\end{equation}
\end{lemma}

\begin{proof}
Put $H_j=X_j(A)-(A-n)Y_j(A)$.  Direct substitution into
\eqref{dc:pell} gives
\[
 H_0=1,\qquad H_{j+1}=nH_j+TY_j(A).
\]
Induction yields
\begin{equation}\label{dc:Hformula}
 H_k=n^k+T\sum_{j=0}^{k-1}n^{k-1-j}Y_j(A).
\end{equation}
Equations \eqref{dc:power} and \eqref{dc:Hformula} show
$m\equiv n^k\pmod T$, independently of the sign of $h$.
Both numbers lie in $[0,T)$, so $m=n^k$.
Substitution and cancellation of $T\ne0$ then give
\eqref{dc:hformula}.  All its summands are nonnegative.
\end{proof}

The exact recurrence already occurs in the inherited positive-power
construction.  The new use is to apply its soundness argument with
$h$ unrestricted and then recover its sign.  The preceding frame and
index analysis verifies every hypothesis of Lemma~\ref{dc:quotientlemma}
for all 23 calls.  We obtain 23 quotient signs; a final elementary pass
obtains the remaining 23 product signs.  All 690 source coordinates
are now nonnegative.

\subsection{Literal interface and absence of circularity}

The semantic interface is checked against the immutable raw gate ledger,
which contains 1,301 coordinates and 1,217 rows.  Source elimination
replaced 611 of these rows by zero identities; the other 606 become
exactly the literal rows of $S$.  The independent checker reconstructs
every omitted raw coordinate by its monic gate definition and verifies
all these polynomial identities.  This reconstruction is valid at
integer zeros, before any nonnegativity assumption on the restored
coordinates.

The interface file then binds \eqref{dc:frame}, \eqref{dc:modulus},
every instance of \eqref{dc:index}, the three addition formulas, and all
23 equations \eqref{dc:power} to exact raw rows and atoms.  Some named
atoms are related by arithmetic gates---for example, both $t$ and
$t+2$ occur in the distinct-atom sum.  Identity comparison preserves
those gate equations; it does not incorrectly declare such atoms
algebraically independent.

The actual inference order, with checked counts, is
\[
\begin{array}{c|c}
\text{Stage}&\text{Source coordinates known nonnegative}\\\hline
131\text{ private guards and elementary propagation}&623\\
\text{frame bounds and further elementary propagation}&644\\
\text{index identification, power recovery, final propagation}&690.
\end{array}
\]
All index premises are checked before invoking the power lemma.
Power outputs enter $W$ as already nonnegative proposed values; no
claim that they are correct is used to derive the frame bounds.
This establishes the sign argument in an ordered form.

\subsection{Integer projection and the resulting theorem}

Let $G_{131}$ be the 606 source rows together with the 131 guards
\eqref{dc:guard}, on 1,083 unrestricted integer coordinates.
The sign proof shows that every integer zero of $G_{131}$ projects
to a natural zero of $S$.  Conversely any natural zero of $S$ extends
by 131 independent three-square representations.  Therefore
\begin{equation}\label{dc:guardedequivalence}
 \exists y\in\N^{690}\ S(y)=0
 \quad\longleftrightarrow\quad
 \exists z\in\Z^{1083}\ G_{131}(z)=0.
\end{equation}

Only after this conversion do we perform signed coordinate elimination.
A row $z_v-f=0$, or its negative, with $f$ an integer polynomial not
involving $z_v$, permits exact substitution everywhere.  Each reduced
integer zero restores uniquely by $z_v=f$; positivity and divisibility
conditions are unnecessary.  The independent checker expands every
substitution and checks that the remaining rows are quadratic.
There are 303 eliminations: 188 affine and 115 quadratic.  No
unconstrained coordinates are discarded.  Reverse traversal restores
the entire guarded solution uniquely, leaving exactly 780 coordinates
and 434 residuals.  Thus
\begin{equation}\label{dc:projectedequivalence}
 \exists z\in\Z^{1083}\ G_{131}(z)=0
 \quad\longleftrightarrow\quad
 \exists x\in\Z^{780}\ P_{780}(x)=0.
\end{equation}
Uniqueness concerns the monic restoration; the initial three-square
extensions need not be unique.

\begin{theorem}\label{dc:main}
Using the inherited calibration \eqref{dc:source}, the explicit
polynomial \eqref{dc:polynomial} satisfies
\[
 \mathrm{ZFC}\vdash
 \bigl(\nexists x\in\Z^{780}:P_{780}(x)=0\bigr)
 \longleftrightarrow\Con(\mathrm{ZFC}).
\]
\end{theorem}
\begin{proof}
Compose \eqref{dc:source}, \eqref{dc:guardedequivalence}, and
\eqref{dc:projectedequivalence}, using that an integer sum of squares
is zero exactly when every summand is zero.  Negate the resulting
existence equivalence.  The new reasoning is internalizable in ZFC
using integer arithmetic, natural induction, standard Pell theory,
and the three-square theorem; it assumes no consistency statement.
\end{proof}

The 36 guards removed from the immediate parent are the 23 private
power quotients, nine moduli $T_n$, the discriminant $D$, and three
coordinates handled by the stronger positive-factor rule:
\texttt{endpoint.ut}, \texttt{factorial0.output}, and
\texttt{crt.nonnegative\_quotient}.  Each guard costs three private
integer coordinates, so the count decreases by $3\cdot36=108$.

\begin{center}
\begin{tabular}{lrr}
\hline
Quantity&Parent&This construction\\\hline
Independent guards&167&131\\
Integer variables&888&780\\
Quadratic residuals&470&434\\
Collected monomials&10,156&9,628\\
Exact degree&4&4\\
Maximum absolute coefficient&$2\,445\,790\,974\,704$&$2\,445\,790\,974\,704$\\
\hline
\end{tabular}
\end{center}
The variable decrease is 12.16\% within this direct quartic family.
Other repository constructions have fewer variables at different
degrees or coefficient costs.  This result makes no claim of global
minimality or worldwide priority.

The independent finite checker imports neither the producer nor its
elimination module.  It verifies source bytes, the staged sign facts,
all semantic interfaces, every signed inverse substitution, and the
complete collected polynomial.  Separate controls require deliberate
coefficient, index, guard and restoration corruptions to be rejected;
the build is reproduced byte for byte in an isolated directory.
These checks authenticate the fixed instance of the proof.  The
unbounded number-theoretic statements and inherited logical calibration
remain the stated mathematical dependencies, and no new end-to-end
Lean or object-ZFC proof object is claimed.


\section{Verification and reproducible artifacts}
\label{sec:verification}

The package is designed so that a reader can inspect an equation without
first running any source program. The ten-positive equation is given by
the full register definitions in this article and by a literal JSON
circuit. The $780$-integer equation is given by an ordered variable list,
all $434$ quadratic residuals, and a separately collected sparse expansion.
The supplied source files and checks define no hidden computation oracle.

\subsection{Ten-positive checks}

The coefficient loader checks both gzip and uncompressed hashes, minimal
binary lengths, bit lengths, $\beta+1=2^{590069}$, and
\eqref{eq:ports}. The builder checks that every circuit input and row is
live, that no forward reference occurs, and that the literal circuit has
exactly $79$ multiplications, $33$ additions, and $21$ subtractions.
Squares and multiplication by a fixed scalar each count as one binary
multiplication. Fixed integer leaves themselves are free in this cost;
the count is not bit complexity.

An independent formula evaluator retains the proof's $W$ abbreviation
while the circuit distributes the corresponding product. Exact modular
comparison checks their agreement at $100$ inputs. Four full univariate
restrictions, over two explicitly recorded prime fields, have degree
$706$ and the proved leading coefficient. The code also checks $150$
finite parity/majorant cases and $172$ finite canonical-gap cases.
The universal upper bound and nonzero top form are mathematical proofs
in the degree section, not conclusions drawn from sampled points.

\subsection{Direct-quartic checks}

The direct checker authenticates the original archive and its selected
members, replays the staged sign certificate against literal rows,
verifies the exact semantic inputs to the shared frame and each power
request, and reconstructs the signed projection. The sparse expansion is
independently collected and compared coefficient for coefficient with
the delivered polynomial. Its nonzero degree-four terms certify exact
degree four. The proof distinguishes these finite algebraic checks from
the inherited unbounded computation calibration.

The top-level replay regenerates outputs in temporary directories and
compares them with the frozen delivery. Checks are run under ordinary
Python and under \texttt{python -O}; correctness conditions use explicit
exceptions rather than relying on assertions that optimization can
disable. A separate AI-assisted review checks the new mathematical proof
ordering. Neither the replay nor the review is described as a proof
assistant checking the complete consistency theorem.

\subsection{Files and exact commands}

\begin{longtable}{@{}>{\raggedright\arraybackslash}p{0.34\textwidth}>{\raggedright\arraybackslash}p{0.61\textwidth}@{}}
\toprule
Path & Purpose\\
\midrule
\endhead
\texttt{report/report.tex}, its two section files, and
\texttt{report/report.pdf} & Editable article and compiled reading copy.\\
\texttt{dyadic10/PROOF.md} & Complete text proof and polynomial recipe.\\
\texttt{dyadic10/circuit.py} and
\texttt{p10\_dyadic\_degree706.json} & All 133 arithmetic rows; ten witnesses;
fixed coefficient binding.\\
\texttt{dyadic10/coefficients/} & The four literal, unchanged coefficient payloads.\\
\texttt{direct780/equations/} & Literal residual equation, full collected polynomial,
sign certificate, signed parent, inverse projection, and exact metrics.\\
\texttt{direct780/PROOF.md} & Complete new delayed-sign and projection proof.\\
\texttt{direct780/upstream/} & Immutable direct parent archive and required extracts.\\
\texttt{provenance/} & Pinned calibration snapshots and source ledger.\\
\texttt{verify.py} and \texttt{VERIFICATION.json} & Whole-package replay and receipt.\\
\texttt{MANIFEST.sha256} & Checksums for the final delivered files.\\
\bottomrule
\end{longtable}

From the extracted package directory, run
\begin{verbatim}
python3 verify.py
python3 -O verify.py
cd report
pdflatex -interaction=nonstopmode -halt-on-error report.tex
pdflatex -interaction=nonstopmode -halt-on-error report.tex
\end{verbatim}
The verification uses the Python standard library. Building the PDF
requires a normal LaTeX installation with the packages named in its
preamble. The mathematical proof can be read without either runtime.

\section{Limitations, interpretation, and further research}
\label{sec:further}

\subsection{What has and has not been achieved}

The concrete outcome is two improvements with explicit domains and
complete polynomial data. The first is a smaller degree at a fixed count
of ten positive unknowns; the second is a smaller integer quartic within
the direct C222 certificate architecture. Both have the required
consistency equivalence after composition with the stated inherited
calibration theorem. The new proofs are ordinary mathematics and use
only finite algebra, divisibility, Pell recurrences, inequalities, and
induction that can be expressed in ZFC.

We have not established that either equation is globally simplest.
The ten-positive equation does not lower the repository's best known
positive variable count from nine. Its $133$-operation circuit is longer
than the parent's $124$-operation circuit. The direct quartic does not
beat the variable counts of the native-coefficient quartics; it retains
the different direct certificate route and its comparatively modest
coefficient height. The parent $407$-natural-variable polynomial is not
claimed as a new result here.

The native tuple and the C222 source calibration are inherited results,
not new independent machine-verification projects. No kernel-checked
endpoint for the whole new equivalence is supplied. We also have not
established priority in the literature for the broad methods of
valuation separation, centering, or recovering quotient signs. The
specific new combinations and their full equation data are the bounded
research claims of this package.

\subsection{Questions suggested by the reductions}

\begin{enumerate}[leftmargin=*]
\item \textbf{Optimize degree without rebuilding root-size barriers.}
The dyadic lemma permits deleting the extra strong scale. Can the
remaining factor $c^2$ in $I=ic^2$ be replaced by a smaller expression
while retaining the exact auxiliary-index theorem? A proposed change
must recover the condition $c\mid m$ used by that theorem; a smaller
formula alone is insufficient.

\item \textbf{Compress the weighted norm evaluator.}
The present circuit uses a general sixteen-fold norm identity.
Can its special weighted and centered arguments be evaluated in fewer
operations without expanding the degree-$706$ polynomial? Any saving
should be verified as equality of the entire fixed polynomial.

\item \textbf{Separate arithmetic and sign inference systematically.}
The $780$-integer construction recovers the signs of power quotients
after identifying their arithmetic values. Which other coordinates have
an analogous delayed sign proof? A useful search should operate on
proof dependencies, not just syntactic nonnegative monomials.

\item \textbf{Optimize the direct system before guarding.}
Which quadratic changes of variables preserve a reverse restoration
and expose more signs without increasing degree? A simultaneous search
over domain proofs and monic eliminations may improve on treating the
two optimizations separately.

\item \textbf{Choose an explicit simplicity objective.}
For example, compare a lexicographic objective of unknown count, degree,
and fixed-payload bytes with a description-length objective charging
every binary numeral and circuit operation. This would distinguish
legitimate improvements from incompatible claims about ``simplest.''

\item \textbf{Close a formalization boundary.}
The new dyadic lemma and the power-quotient recurrence are small enough
to formalize separately. A full endpoint additionally needs the fixed
polynomial semantics, source calibration, and derivation inside the
chosen formal presentation of ZFC. These are separate obligations from
testing a builder.
\end{enumerate}

\appendix
\section{Domain conventions and conversion}

Throughout, $\N=\{0,1,2,\ldots\}$ and $\Pos=\{1,2,3,\ldots\}$.
The ten-positive construction can be viewed as a ten-natural equation
by substituting $x_j+1$ for each unknown. This preserves its exact degree
and adds at most ten additions to the displayed circuit.
It does not make the unknowns unrestricted integers.

If an elementary all-integer conversion is wanted for that same
polynomial, the identity
\[
 u>0\quad\longleftrightarrow\quad
 \exists a,b,c\in\Z\;u=1+a^2+b^2+c(c+1)
\]
follows from the three-squares theorem (equivalently the standard
representation of nonnegative integers used in \cite{sun2026}).
Applying it independently gives a polynomial in $30$ unrestricted
integers of degree $1412$. This is merely an explicit domain conversion,
not a competitive new endpoint, so it is not counted as a third result.

\section{Attribution and source access}

The project snapshot is
\begin{center}\small\sha.\end{center}
The native coefficient files come from package 44; the calibrated source
interface is package 61; the immediate positive parent is the in-repository
degree-$816$ candidate. That candidate acknowledges earlier four-root,
centering, and residue constructions. This work preserves those credits
and adds its own concrete dyadic descaling and degree computation.
The direct source is the immutable incoming $407/888$ archive. Its
shared-frame and direct certificate results are retained as inherited
material. The new sign recovery and the resulting $780$-variable data
are distinguished from that archive.

The package license is MIT No Attribution for new material. Inherited
licenses and notices are retained with the material to which they apply.
All implementation and review in the present package are AI-assisted;
the independent review means a separate analysis task, not an external
journal referee or a proof assistant.

\begin{thebibliography}{9}
\small\raggedright
\bibitem{parent816}
ProveIt contributors, \emph{Ten positive unknowns, degree 816: a scaled,
centered sixteen-fold norm}, candidate and proof at snapshot \sha.
\href{https://github.com/VladimirReshetnikov/ProveIt/blob/ec80841c07e91be8956ce6ed4af9a9c0b782d3d9/Computability/DiophantineConsistency/candidates/ten-positive-scaled-norm/PROOF.md}
{Pinned complete parent proof}. Included unchanged in
\texttt{dyadic10/upstream/PROOF816.md}.

\bibitem{source61}
ProveIt contributors, \emph{Congruence carriers}, package 61,
Theorem \texttt{thm:source}, canonical promises, and Pell lemmas.
\href{https://github.com/VladimirReshetnikov/ProveIt/blob/ec80841c07e91be8956ce6ed4af9a9c0b782d3d9/Computability/DiophantineConsistency/con-zfc-diophantine-forms/61-congruence-carriers/report/congruence_carriers.tex}
{Pinned LaTeX source}.

\bibitem{source44}
ProveIt contributors, \emph{Complement-packed native clauses: a smaller
complete coefficient tuple}, package 44, native proof and binding.
\href{https://github.com/VladimirReshetnikov/ProveIt/blob/ec80841c07e91be8956ce6ed4af9a9c0b782d3d9/Computability/DiophantineConsistency/con-zfc-diophantine-forms/44-b4-sparse/native/PROOF.md}
{Pinned native proof}.

\bibitem{direct407}
ProveIt incoming research package, \emph{ConZFC\_407\_888\_Quartics},
archive SHA-256
\texttt{59f847872298260ee35e9d8e4cd0b52d2d78658ac523981fe060e81f40ca4a59}.
\href{https://github.com/VladimirReshetnikov/ProveIt/blob/ec80841c07e91be8956ce6ed4af9a9c0b782d3d9/docs/incoming/ConZFC_407_888_Quartics_2026-10-24.zip}
{Pinned original archive}. Included unchanged under \texttt{direct780/upstream/}.

\bibitem{bayer2026}
J. Bayer, M. David, M. Ha\ss ler, Y. Matiyasevich, and D. Schleicher,
\emph{Diophantine equations over $\Z$: universal bounds and parallel
formalization}, arXiv:2506.20909v3, 22 October 2026.
\url{https://arxiv.org/html/2506.20909v3}.

\bibitem{sun2026}
Z.-W. Sun, \emph{Ten unknowns for Hilbert's tenth problem over the
integers}, arXiv:2609.01594v2, 2 September 2026.
\url{https://arxiv.org/html/2609.01594v2}.
\end{thebibliography}
\end{document}
